% Conceptual: serial AR decode vs draft--verify. Caption lives in thesis.tex.
\begin{tikzpicture}[
  font=\small,
  >=Stealth,
  pass/.style={
    draw=tgtC,
    fill=tgtC!12,
    rounded corners=2pt,
    minimum width=2.2cm,
    minimum height=0.72cm,
    align=center,
    inner sep=2pt,
    font=\footnotesize,
  },
  dpass/.style={
    draw=dftC,
    fill=dftC!12,
    rounded corners=2pt,
    minimum width=1.7cm,
    minimum height=0.58cm,
    align=center,
    inner sep=1.5pt,
    font=\footnotesize,
  },
  tok/.style={
    draw,
    rounded corners=1.5pt,
    minimum width=0.72cm,
    minimum height=0.52cm,
    inner sep=1pt,
    font=\ttfamily\footnotesize,
  },
]
\node[anchor=west, font=\bfseries\small] at (0,4.85) {%
  (α)~Αυτοπαλινδρομικό decode: ένα πέρασμα, ένα token};

\node[pass] (P0) at (1.20,3.95) {στόχος};
\node[tok, fill=tgtC!20] (T0) at (2.90,3.95) {$t_1$};
\draw[->, thick, tgtC] (P0) -- (T0);

\node[pass] (P1) at (4.85,3.95) {στόχος};
\node[tok, fill=tgtC!20] (T1) at (6.55,3.95) {$t_2$};
\draw[->, thick, tgtC] (P1) -- (T1);

\node[pass] (P2) at (8.50,3.95) {στόχος};
\node[tok, fill=tgtC!20] (T2) at (10.20,3.95) {$t_3$};
\draw[->, thick, tgtC] (P2) -- (T2);
\node[font=\footnotesize, text=arC] at (11.55,3.95) {$\cdots$};

\draw[->, thick, arC] (T0) -- (P1);
\draw[->, thick, arC] (T1) -- (P2);

\node[anchor=west, font=\footnotesize, text=arC, text width=13.5cm] at (0,3.15) {%
  Κάθε βέλος είναι πλήρες forward pass: σχεδόν όλα τα βάρη διαβάζονται για
  \emph{ένα} νέο token. Σε μικρό batch η GPU περιμένει τη μνήμη.};

\node[anchor=west, font=\bfseries\small] at (0,2.35) {%
  (β)~Κερδοσκοπία: φθηνή πρόταση, μία batched επαλήθευση};

\node[dpass] (D0) at (0.95,1.45) {drafter};
\node[tok, fill=dftC!25, draw=dftC] (D1) at (2.40,1.45) {$d_1$};
\node[tok, fill=dftC!25, draw=dftC] (D2) at (3.30,1.45) {$d_2$};
\node[tok, fill=dftC!25, draw=dftC] (D3) at (4.20,1.45) {$d_3$};
\node[tok, fill=dftC!25, draw=dftC] (D4) at (5.10,1.45) {$d_4$};
\draw[->, dftC] (D0) -- (D1);
\draw[->, dftC] (D1) -- (D2);
\draw[->, dftC] (D2) -- (D3);
\draw[->, dftC] (D3) -- (D4);

\node[pass, minimum width=4.35cm] (V) at (9.15,1.45) {στόχος: ένα batched πέρασμα};
\draw[->, ultra thick, tgtC] (D4.east) -- (V.west);

\node[tok, fill=dftC!35, draw=dftC] (A1) at (2.40,0.40) {$d_1$};
\node[font=\scriptsize, text=dftC] at (2.40,0.82) {$\checkmark$};
\node[tok, fill=dftC!35, draw=dftC] (A2) at (3.30,0.40) {$d_2$};
\node[font=\scriptsize, text=dftC] at (3.30,0.82) {$\checkmark$};
\node[tok, fill=rejC!20, draw=rejC] (X) at (4.20,0.40) {$x$};
\node[font=\scriptsize, text=rejC] at (4.20,0.82) {$\times$};
\node[tok, dashed, draw=arC, text=arC] (R) at (5.10,0.40) {$d_4$};

\node[anchor=west, font=\footnotesize, text=inkC, text width=13.5cm] at (0,-0.35) {%
  Αποδοχή $d_1,d_2$. Απόρριψη του $d_3$: ο στόχος εκπέμπει διόρθωση $x$.
  Το $d_4$ απορρίπτεται χωρίς να χρησιμοποιηθεί.};

\node[anchor=west, font=\footnotesize, text=arC, text width=13.5cm] at (0,-1.05) {%
  Ο στόχος δεν παρακάμπτεται και δεν προσεγγίζεται. Τρέχει σε κάθε γύρο·
  η ωφέλεια είναι περισσότερα εκπεμπόμενα tokens ανά ανάγνωση των βαρών του.};
\end{tikzpicture}
